\section{Optimal Matching between Point Clouds}
\label{sec-matching}
\subsection{Monge Problem for Discrete Points}
\label{sec-monge-pbm}
\paragraph{Assignment problem.}
Let $\C\in\RR^{n\times n}$ be a cost matrix, where $\C_{i,j}$ is the cost of pairing source $i$ with target $j$, and let $\Perm(n)$ denote the set of bijections of $\range{n}=\{1,\ldots,n\}$. The optimal assignment problem is

\eql{\label{eq-optimal-assignment}
\umin{\sigma \in \Perm(n)} \frac{1}{n}\sum_{i=1}^n \C_{i,\sigma(i)}.
}
\paragraph{Convex costs on the line.}
\begin{proposition}[Monotone matching on the line]\label{prop-matching-1d-monotone}
Assume that $x_1,\ldots,x_n$ are pairwise distinct and that $y_1,\ldots,y_n$ are pairwise distinct. If $\C_{i,j}=h(x_i-y_j)$ for a strictly convex function $h:\RR\to\RR$, then the unique optimizer is the order-preserving permutation, characterized by
\((x_i-x_{i'})(y_{\sigma(i)}-y_{\sigma(i')})>0 \qquad\text{for every }i\neq i'.\)
In particular, the result applies to $h(s)=|s|^p$ for $p>1$.
\end{proposition}
\begin{proof}
A permutation that does not preserve order contains an inversion: after relabeling, $x<x'$ are matched to $y'>y$. Set $d=y'-y>0$ and
\(D(s)=\frac{h(s)-h(s-d)}{d}.\)
Strict convexity of $h$ makes $D$ strictly increasing. Therefore
\[
\begin{split}
&h(x-y')+h(x'-y)-h(x-y)-h(x'-y')\\
&\hspace{8em}=d\bigl(D(x'-y)-D(x-y)\bigr)>0.
\end{split}
\]
Swapping the two inverted targets strictly lowers the cost. Repeating this exchange eliminates every inversion; the only order-preserving bijection between the two sorted clouds pairs equal ranks.
\end{proof}
If $h$ is convex but not strictly convex, the same exchange inequality is non-strict. The equal-rank matching remains optimal, but other, non-monotone optimizers can coexist. To construct a monotone optimizer, choose sorting permutations $\sigma_X,\sigma_Y$ such that

\(x_{\sigma_X(1)} \leq x_{\sigma_X(2)} \leq \ldots \qandq y_{\sigma_Y(1)} \leq y_{\sigma_Y(2)} \leq \ldots\)
and match $x_{\sigma_X(k)}$ with $y_{\sigma_Y(k)}$. Thus $\sigma=\sigma_Y\circ\sigma_X^{-1}$ is optimal.

\paragraph{Concave costs on the line.}
\(p_1<q_1<p_2<q_2<\cdots<p_N<q_N,\)
with the opposite orientation handled by exchanging the roles of $p$ and $q$, Delon, Salomon and Sobolevski define the local indicators

\begin{align*}
I_k^p(i)
&=
c(p_i,q_{i+k})
+
\sum_{r=0}^{k-1}c(p_{i+r+1},q_{i+r})
-
\sum_{r=0}^{k}c(p_{i+r},q_{i+r}),\\
I_k^q(i)
&=
c(p_{i+k+1},q_i)
+
\sum_{r=1}^{k}c(p_{i+r},q_{i+r})
-
\sum_{r=0}^{k}c(p_{i+r+1},q_{i+r}),
\end{align*}
For arbitrary real masses, the stratification argument in the same work gives a larger $O(n^3)$ worst-case bound. Very concave costs also motivate the simpler greedy rule that repeatedly matches the closest red--blue pair. These methods are not generic linear-programming solvers: they exploit the order and concavity specific to the line.

\paragraph{Histogram equalization.}
\paragraph{Flat directions for the linear cost.}
\paragraph{Optimal transport on the circle.}\label{rem-circle-ot-cut}
\(d_{\mathbb S^1}(x,y):=\min_{k\in\ZZ}|x-y+k|, \qquad c_p(x,y):=d_{\mathbb S^1}(x,y)^p,\qquad p>1.\)
\begin{proposition}[Discrete circle transport by a cut]\label{prop-circle-ot-cut}
Assume that the $2n$ points $x_1,\ldots,x_n,y_1,\ldots,y_n$ are pairwise distinct. Let $x_{(1)},\ldots,x_{(n)}$ and $y_{(1)},\ldots,y_{(n)}$ be fixed cyclic orderings, with indices understood modulo $n$. For $p>1$, every optimal assignment for the cost $c_p$ is a cyclic shift
\(x_{(k)}\longmapsto y_{(k+s)}, \qquad k\in\range{n},\quad s\in\{0,\ldots,n-1\}.\)
\[
E_s
\eqdef
\frac1n\sum_{k=1}^n
d_{\mathbb S^1}\!\left(x_{(k)},y_{(k+s)}\right)^p
\]
over the $n$ shifts. For each optimal assignment there is a cut $\theta\in\mathbb S^1\setminus(\{x_i\}_i\cup\{y_j\}_j)$, crossed by none of its shortest transport arcs, such that the lifted assignment on $(\theta,\theta+1)$ is the equal-rank matching on the line.
\end{proposition}
\begin{proof}
Fix an optimal assignment $\sigma$ and let $\gamma_i$ be an open shortest arc from $x_i$ to $y_{\sigma(i)}$, with either orientation fixed in the antipodal case. The path-uncrossing lemma of Delon, Rabin and Gousseau states that if two arcs of an optimal assignment intersect, then they have the same orientation and neither is strictly contained in the other. Its proof is a two-edge exchange: opposite orientations or strict containment would, by monotonicity and strict convexity of $r\mapsto r^p$, produce a strictly cheaper pairing of the four endpoints.

Suppose that the arcs cover the circle. Since each $\gamma_i$ is open, every source point lies in another transport arc. Order the sources cyclically. The path-uncrossing lemma then propagates one common orientation: either $x_{(k+1)}\in\gamma_{(k)}$ for every $k$, or $x_{(k-1)}\in\gamma_{(k)}$ for every $k$. In the first case, define a new assignment by giving $x_{(k+1)}$ the target previously assigned to $x_{(k)}$; every transport arc becomes strictly shorter. The opposite orientation gives the analogous backward reassignment. Both contradict optimality. Hence some $\theta$ lies outside every transport arc.

Cut the circle at such a $\theta$ and lift it to $(\theta,\theta+1)$. Every matched geodesic avoids the cut, so its circular length equals the ordinary distance between the corresponding lifts. Proposition~\ref{prop-matching-1d-monotone} now implies that the lifted matching preserves order. Returning to the circle, an order-preserving bijection between two cyclically ordered $n$-point sets is a cyclic shift. Minimizing over these $n$ shifts therefore recovers the optimum.
\end{proof}
\paragraph{Rational weights.}
\begin{proposition}[Rational weights as duplicated uniform matching]\label{prop-rational-weights-duplicated-matching}
Let
\(\alpha=\sum_{i=1}^n \frac{k_i}{N}\delta_{x_i}, \qquad \beta=\sum_{j=1}^m \frac{\ell_j}{N}\delta_{y_j}, \qquad \sum_i k_i=\sum_j\ell_j=N,\)
with positive integers $k_i,\ell_j$. Replace each $x_i$ by $k_i$ identical copies and each $y_j$ by $\ell_j$ identical copies, producing two uniform $N$-point clouds. Then the duplicated assignment problem and the discrete Kantorovich problem between $\alpha$ and $\beta$ have the same optimal value. Moreover, there exists an optimal coupling of the form
\(P_{ij}=\frac{n_{ij}}{N}, \qquad n_{ij}\in\NN, \qquad \sum_j n_{ij}=k_i, \quad \sum_i n_{ij}=\ell_j,\)
and couplings whose scaled entries $NP_{ij}$ are integral are exactly the collapsed assignments between duplicated clouds. Fractional optimal couplings may nevertheless coexist when the optimum is degenerate.
\end{proposition}
\begin{proof}
Any assignment between the duplicated source and target clouds defines integers $n_{ij}$ counting how many copied particles of type $x_i$ are matched to copied particles of type $y_j$. These counts satisfy $\sum_j n_{ij}=k_i$ and $\sum_i n_{ij}=\ell_j$, and the associated coupling $P_{ij}=n_{ij}/N$ has marginals $k_i/N$ and $\ell_j/N$. The assignment cost is
\(\frac1N\sum_{i,j} n_{ij}c(x_i,y_j) = \sum_{i,j}P_{ij}c(x_i,y_j).\)
Conversely, any nonnegative integer count matrix with those row and column sums can be realized by allocating the $k_i$ copies of each $x_i$ among the target copies according to $(n_{ij})_j$. Scale a feasible coupling by $N$ and write $Q=NP$. The constraints on $Q$ have integer right-hand sides $(k_1,\ldots,k_n,\ell_1,\ldots,\ell_m)$. After multiplying the target rows by $-1$, their coefficient matrix is the oriented node-edge incidence matrix of a bipartite graph and is therefore totally unimodular. Every vertex of this transportation polytope is integral. Since a linear objective attains its minimum at a vertex, an integral optimal $Q$ exists, proving equality of the two optimal values. The argument establishes existence, not integrality of every optimizer; convex combinations of distinct integral optima give fractional optima.
\end{proof}
\paragraph{Two-dimensional geometry.}
\begin{proposition}[No proper crossings in Euclidean-distance matchings]\label{prop-planar-no-proper-crossing}
Let $x_1,\ldots,x_n,y_1,\ldots,y_n\in\RR^2$ and let $c(x,y)=\norm{x-y}$. If $\sigma$ is optimal, then two matched segments cannot cross properly: their relative interiors cannot meet at a point where their supporting lines are distinct.
\end{proposition}
\begin{proof}
Suppose that $[x_i,y_{\sigma(i)}]$ and $[x_j,y_{\sigma(j)}]$ cross properly at $z$. The triangle inequality along the two reconnected paths gives \(\norm{x_i-y_{\sigma(j)}} \leq \norm{x_i-z}+\norm{z-y_{\sigma(j)}},\)
and
\(\norm{x_j-y_{\sigma(i)}} \leq \norm{x_j-z}+\norm{z-y_{\sigma(i)}}.\)
The sum of the two right-hand sides is exactly the cost of the two original segments. At least one inequality is strict because a proper crossing is non-collinear. Swapping $\sigma(i)$ and $\sigma(j)$ therefore strictly decreases the assignment cost, contradicting optimality. Collinear overlaps are excluded from the statement precisely because the exchange can then have equal cost.
\end{proof}
\(C_n=\frac{1}{n+1}\binom{2n}{n}\sim \frac{4^n}{\sqrt{\pi}n^{3/2}}.\)
\subsection{Matching Algorithms}
\paragraph{Classical assignment methods.}
Efficient algorithms avoid enumerating $n!$ permutations. Auction algorithms attach prices to targets: an unmatched source bids for the target of smallest reduced cost, equivalently largest reduced profit, and raises that target's price.

\paragraph{Hungarian primal-dual method.}
\(u_i+v_j\leq \C_{i,j}\qquad\forall i,j.\)
\[
\delta=\min_{i\in S,\ j\notin T}\bigl(\C_{i,j}-u_i-v_j\bigr),
\qquad
u_i\leftarrow u_i+\delta\ (i\in S),\qquad
v_j\leftarrow v_j-\delta\ (j\in T).
\]
\begin{prop}[Correctness and complexity of the Hungarian primal-dual method]\label{prop-hungarian-correct}
Assume the Hungarian method terminates with a perfect matching $\sigma$ contained in the equality graph
\(E(u,v)=\enscond{(i,j)}{u_i+v_j=\C_{i,j}},\)
where $(u,v)$ is dual feasible, i.e. $u_i+v_j\leq \C_{i,j}$ for all $(i,j)$. Then $\sigma$ is an optimal assignment. Moreover, the usual Hungarian updates terminate after finitely many augmentations.
With maintained slacks, the method uses $O(n^3)$ arithmetic operations.
\end{prop}
\begin{proof}
For any permutation $\tau$, dual feasibility gives \(\sum_i \C_{i,\tau(i)} \geq \sum_i (u_i+v_{\tau(i)}) = \sum_i u_i+\sum_j v_j.\)
This is the weak duality lower bound. If $\sigma$ is contained in the equality graph, then
\(\sum_i \C_{i,\sigma(i)} = \sum_i u_i+\sum_j v_j,\)
so the primal cost of $\sigma$ reaches the dual lower bound and is optimal.

At each successful augmentation, the matching cardinality increases by one, so there are exactly $n$ augmentation phases. During one phase, the algorithm grows an alternating tree in the equality graph. Before reaching a free target, the tree invariant is $|S|=|T|+1$, so $T$ cannot contain every target. If $N_E(S)\setminus T$ is empty, every slack from $S$ to $T^c$ is positive and the minimum defining $\delta$ exists. For edges in $S\times T$, adding $\delta$ to source labels and subtracting it from target labels preserves tightness. For edges in $S\times T^c$, the definition of $\delta$ preserves feasibility and makes at least one new edge tight; edges in $S^c\times T$ become looser, and all remaining inequalities are unchanged. Thus the reachable sets strictly grow after every dual update and can grow at most $n$ times in one phase. If the slacks $\min_{i\in S}(\C_{i,j}-u_i-v_j)$ are maintained when a source enters $S$, each tree expansion costs $O(n)$. A phase costs $O(n^2)$ and the $n$ phases cost $O(n^3)$. Hence the method reaches a perfect optimal matching.
\end{proof}